\documentclass[12pt]{article}

% -------------------- Idioma y codificación --------------------
\usepackage[spanish]{babel}
\addto\captionsspanish{\renewcommand{\tablename}{Tabla}}
\usepackage[utf8]{inputenc}
\usepackage[T1]{fontenc}

% -------------------- Formato de página --------------------
\usepackage{geometry}
\geometry{margin=2.5cm}
\usepackage{setspace}
\setstretch{1.15}

% -------------------- Colores --------------------
\usepackage{xcolor}
\definecolor{tituloazul}{RGB}{31,100,160}
\definecolor{barraazul}{RGB}{210,230,250}

% -------------------- Links clickeables --------------------
\usepackage[
    colorlinks=true,
    linkcolor=tituloazul,
    citecolor=tituloazul,
    urlcolor=tituloazul
]{hyperref}

% -------------------- Estilo de secciones --------------------
\usepackage{titlesec}

\titleformat{\section}
{\color{tituloazul}\large\bfseries}
{\thesection}{1em}{}

\titleformat{\subsection}
{\color{tituloazul}\normalsize\bfseries}
{\thesubsection}{1em}{}

\titleformat{\subsubsection}
{\color{tituloazul}\small\bfseries}
{\thesubsubsection}{0.5em}{}

\titlespacing{\section}{0pt}{12pt}{6pt}
\titlespacing{\subsection}{0pt}{8pt}{4pt}

% -------------------- Encabezados --------------------
\usepackage{fancyhdr}
\pagestyle{fancy}

\fancyhf{}
\fancyhead[L]{\textcolor{tituloazul}{Electrónica de Potencia}}
\fancyhead[R]{\textcolor{tituloazul}{Práctica 2}}
\fancyfoot[C]{\thepage}

\setlength{\headheight}{15pt}

% -------------------- Figuras --------------------
\usepackage{graphicx}
\usepackage{caption}
\usepackage{subcaption}
\usepackage{float}

\captionsetup{
    skip=5pt,
    font=small
}

% ============================================================
% IMPORTANTE:
% Las imágenes deben estar dentro de la carpeta "figuras/"
% junto al archivo .tex.
%
% Ejemplo:
%
% informe.tex
% figuras/
%   90 grados salida lm358 a0.jfif
%   120 grados salida lm358 a0.jfif
%   ...
%
% Se utiliza \detokenize{} para permitir espacios en los nombres.
% ============================================================

% -------------------- Matemáticas --------------------
\usepackage{amsmath}

\usepackage[
    spanish,
    capitalize,
    nameinlink
]{cleveref}

\crefname{figure}{figura}{figuras}
\Crefname{figure}{Figura}{Figuras}

\crefname{equation}{ecuación}{ecuaciones}
\Crefname{equation}{Ecuación}{Ecuaciones}

\crefname{table}{tabla}{tablas}
\Crefname{table}{Tabla}{Tablas}

\usepackage{siunitx}

\sisetup{
    per-mode=symbol,
    detect-all
}

% -------------------- Tablas --------------------
\usepackage{booktabs}
\usepackage{enumitem}

\setlist{
    topsep=5pt,
    itemsep=3pt,
    parsep=2pt
}

% -------------------- Circuitos --------------------
\usepackage[siunitx]{circuitikz}
\usetikzlibrary{babel}

% -------------------- Referencias --------------------
\usepackage{cite}

% -------------------- Estudiantes --------------------
\newcommand{\estudiantes}{
Jose Leonardo Perez Poveda, Sergio Hoyos Mejía, Alberto Mario Salazar Callejas
}

% ============================================================
% DOCUMENTO
% ============================================================

\begin{document}

% -------------------- Título --------------------
\vspace*{-1cm}

\begin{center}

{\Huge
\textcolor{tituloazul}{
\textbf{Práctica 2}\\
Interruptores de Electrónica de Potencia y Circuitos de Disparo. Parte 2
}}

\vspace{0.2cm}

{\textbf{\estudiantes}}

\vspace{0.2cm}

\textit{Ingeniería Electrónica}\\
\textit{Universidad de Antioquia}\\
\textit{Medellín-Colombia}

\end{center}

\vspace{0.3cm}

% ============================================================
% RESUMEN
% ============================================================

\section{Informe}

La práctica permitió implementar y analizar un sistema de control de potencia por ángulo de fase utilizando un SCR BT151-500R, un puente rectificador y un optotriac MOC3020. El sistema sincroniza el disparo del tiristor con la red eléctrica mediante una etapa de detección de cruce por cero aislada.

El microcontrolador recibe la señal de sincronismo proveniente del RH111 y genera un pulso de disparo después de un tiempo determinado por el ángulo $\alpha$. Este pulso activa el transistor Q1 y posteriormente el MOC3020, que proporciona aislamiento galvánico antes de llevar la señal hacia la compuerta del SCR.

Adicionalmente, la tensión sobre la carga es medida mediante un divisor resistivo, un LM358 y el optoacoplador lineal IL300. Esta etapa permite enviar información analógica de la carga hacia el microcontrolador sin establecer una conexión eléctrica directa entre la etapa de potencia y la electrónica de control.

% ============================================================
% 1. CONTROL Y DISPARO
% ============================================================

\subsection{Señal de control y tensión en la compuerta del SCR}

La salida del microcontrolador genera un pulso de aproximadamente $5\,V$, cuya posición dentro del ciclo de la red depende del ángulo de disparo $\alpha$.

La cadena de activación es:

\[
\text{Microcontrolador}
\rightarrow R4
\rightarrow Q1
\rightarrow MOC3020
\rightarrow R5
\rightarrow Gate
\]

El transistor Q1, correspondiente a un 2N2222, funciona como interruptor y etapa de excitación para el LED interno del MOC3020. De esta manera, el microcontrolador no necesita suministrar directamente toda la corriente requerida por el optotriac.

El MOC3020 proporciona el aislamiento galvánico entre la etapa de control y la etapa de potencia. Su característica de disparo aleatorio permite activar el SCR en el instante determinado por el ángulo $\alpha$.

La señal observada en el Gate corresponde finalmente al pulso que permite que el SCR pase del estado de bloqueo al estado de conducción.

% ============================================================
% FIGURA: 90 GRADOS RH111 + DISPARO
% ============================================================

\begin{figure}[H]
    \centering
    \begin{subfigure}{0.45\textwidth}
        \centering
        \includegraphics[
            width=\textwidth
        ]{\detokenize{figuras/señal rectificada y disparo a 90.png}}
        \caption{Señal de entrada del RH111 $\alpha=90^\circ$ y disparo del microcontrolador.}
    \end{subfigure}
    \hfill
    \begin{subfigure}{0.45\textwidth}
        \centering
        \includegraphics[
            width=\textwidth
        ]{\detokenize{lab2/figuras/señal rectificada y disparo a 120.png}}
        \caption{Señal para $\alpha=120^\circ$.}
    \end{subfigure}
    \hfill
    \begin{subfigure}{0.45\textwidth}
        \centering
        \includegraphics[
            width=\textwidth
        ]{\detokenize{lab2/figuras/señal rectificada y disparo a 170.png}}
        \caption{Señal para $\alpha=170^\circ$.}
    \end{subfigure}

    \caption{Sincronización y señal de disparo para diferentes ángulos.}
    \label{fig:disparo_90_120}
\end{figure}

% ============================================================
% 2. CONTROL DE FASE
% ============================================================

\subsection{Forma de onda de la tensión sobre la carga}

Para una red de aproximadamente $120\,V_{\mathrm{RMS}}$ y $60\,Hz$, la tensión pico es:

\[
V_{\mathrm{pico}}
=
120\sqrt{2}
\approx
170\,V
\]

y el período de la señal es:

\[
T=\frac{1}{60}
\approx16.67\,ms
\]

El tiempo correspondiente a un determinado ángulo de disparo puede calcularse mediante:

\[
t_{\alpha}
=
\frac{\alpha}{360^\circ}T
\]

Por ejemplo, para $\alpha=60^\circ$:

\[
t_{\alpha}
=
\frac{60^\circ}{360^\circ}
(16.67\,ms)
\approx
2.78\,ms
\]

Durante el intervalo previo al disparo el SCR permanece bloqueado, por lo que la tensión sobre la carga es prácticamente nula. Después del disparo, el SCR conduce y la tensión de la red aparece sobre la carga.

Por lo tanto:

\[
\alpha \uparrow
\quad\Rightarrow\quad
V_{\mathrm{carga,RMS}}\downarrow
\]

mientras que:

\[
\alpha \downarrow
\quad\Rightarrow\quad
V_{\mathrm{carga,RMS}}\uparrow
\]

% ============================================================
% FIGURAS BOMBILLO
% ============================================================

\begin{figure}[H]
    \centering

    \begin{subfigure}{0.31\textwidth}
        \centering
        \includegraphics[
            width=\textwidth
        ]{\detokenize{lab2/figuras/20 paralelo bombillo.png}}
        \caption{$\alpha=20^\circ$.}
    \end{subfigure}
    \hfill
    \begin{subfigure}{0.31\textwidth}
        \centering
        \includegraphics[
            width=\textwidth
        ]{\detokenize{lab2/figuras/90 paralelo bombillo.png}}
        \caption{$\alpha=90^\circ$.}
    \end{subfigure}
    \hfill
    \begin{subfigure}{0.31\textwidth}
        \centering
        \includegraphics[
            width=\textwidth
        ]{\detokenize{lab2/figuras/120 paralelo bombillo.png}}
        \caption{$\alpha=120^\circ$.}
    \end{subfigure}

    \caption{Tensión observada en paralelo con el bombillo para diferentes ángulos de disparo.}
    \label{fig:bombillo_angulo}
\end{figure}

La comparación de las formas de onda evidencia el principio de control por fase. Para un ángulo pequeño, el SCR conduce durante una mayor parte de cada semiciclo, mientras que al aumentar $\alpha$ disminuye el intervalo de conducción.

La tensión sobre el SCR presenta el comportamiento complementario: antes del disparo soporta gran parte de la tensión de la etapa de potencia y, una vez que entra en conducción, su tensión ánodo-cátodo disminuye hasta un valor cercano a su caída directa.

% ============================================================
% 3. IL300
% ============================================================

\subsection{Medición aislada mediante el IL300}

El IL300 cumple una función diferente a la del MOC3020. Mientras el MOC3020 transmite la orden de disparo desde el microcontrolador hacia la etapa de potencia, el IL300 permite transferir información analógica desde la etapa de potencia hacia el microcontrolador.

La cadena de medición es:

\[
V_{\mathrm{carga}}
\rightarrow
\text{Divisor resistivo}
\rightarrow
IL300
\rightarrow
LM358
\rightarrow
A0
\]

El divisor resistivo reduce la tensión de la carga a un nivel adecuado para la etapa de medición. Posteriormente, el IL300 realiza la transferencia óptica manteniendo el aislamiento galvánico. Finalmente, el LM358 acondiciona la señal para la entrada analógica A0.

El valor medido por el microcontrolador depende de la tensión real sobre la carga, del divisor resistivo, de la transferencia del IL300 y de la ganancia del amplificador operacional.

% ============================================================
% FIGURAS A0
% ============================================================

\begin{figure}[H]
    \centering

    \begin{subfigure}{0.31\textwidth}
        \centering
        \includegraphics[
            width=\textwidth
        ]{\detokenize{lab2/figuras/90 grados salida lm358 a0.png}}
        \caption{$\alpha=90^\circ$.}
    \end{subfigure}
    \hfill
    \begin{subfigure}{0.31\textwidth}
        \centering
        \includegraphics[
            width=\textwidth
        ]{\detokenize{lab2/figuras/120 grados salida lm358 a0.png}}
        \caption{$\alpha=120^\circ$.}
    \end{subfigure}
    \hfill
    \begin{subfigure}{0.31\textwidth}
        \centering
        \includegraphics[
            width=\textwidth
        ]{\detokenize{lab2/figuras/170 grados salida lm358 a0.png}}
        \caption{$\alpha=170^\circ$.}
    \end{subfigure}

    \caption{Señal acondicionada por el LM358 observada en la entrada A0.}
    \label{fig:a0_angulos}
\end{figure}

Las mediciones muestran que la señal entregada a A0 permite identificar la forma de onda asociada a la tensión de carga. Al aumentar el ángulo de disparo, se reduce el intervalo durante el cual existe tensión sobre el bombillo, por lo que la señal medida presenta un menor intervalo de conducción.

% ============================================================
% 4. DIFERENCIA MOC3020 IL300
% ============================================================

\subsection{Diferencia entre el MOC3020 y el IL300}

\begin{table}[H]
\centering
\caption{Comparación entre los optoacopladores utilizados.}
\begin{tabular}{p{3.2cm} p{5.2cm} p{5.2cm}}
\toprule
\textbf{Característica} &
\textbf{MOC3020} &
\textbf{IL300} \\
\midrule
Tipo &
Optotriac &
Optoacoplador lineal \\
\midrule
Función &
Disparo &
Medición y realimentación \\
\midrule
Señal &
Pulso de control &
Señal analógica \\
\midrule
Dirección &
Control $\rightarrow$ potencia &
Potencia $\rightarrow$ control \\
\midrule
Aplicación &
Activación del SCR &
Medición aislada de la carga \\
\bottomrule
\end{tabular}
\end{table}

En resumen, el MOC3020 transmite una orden de conmutación, mientras que el IL300 transmite información analógica. Ambos mantienen aislamiento galvánico, pero sus funciones dentro del sistema son diferentes.

% ============================================================
% 5. VARIACIÓN DEL ÁNGULO
% ============================================================

\subsection{Variación del ángulo de disparo}

La modificación de $\alpha$ desde el microcontrolador cambia directamente el instante en que el SCR comienza a conducir.

Para una carga resistiva, un ángulo pequeño produce mayor tiempo de conducción y, por tanto, mayor potencia entregada. Al aumentar el ángulo, el tiempo de conducción disminuye y también lo hacen la tensión RMS y la potencia de la carga.

La relación observada experimentalmente puede resumirse como:

\[
\alpha\downarrow
\Rightarrow
\text{mayor conducción}
\Rightarrow
\text{mayor potencia}
\]

\[
\alpha\uparrow
\Rightarrow
\text{menor conducción}
\Rightarrow
\text{menor potencia}
\]

Los resultados obtenidos con el bombillo confirman esta tendencia, especialmente al comparar las señales correspondientes a $20^\circ$, $90^\circ$ y $120^\circ$.

% ============================================================
% 6. PULSO DE DISPARO
% ============================================================

\subsection{Efecto de un pulso demasiado corto}

El SCR requiere que la corriente de Gate sea suficiente para superar la corriente de disparo $I_{GT}$ durante un tiempo adecuado.

Si el pulso transmitido mediante el MOC3020 es demasiado corto, pueden producirse disparos intermitentes o ausencia de conducción. Esto genera una forma de onda irregular en la carga y variaciones de potencia entre ciclos.

Para un disparo confiable debe cumplirse: $I_G \geq I_{GT}$


durante un intervalo suficiente para que la corriente principal del SCR alcance la corriente de enclavamiento $I_L$.

Una vez enclavado, el SCR permanece conduciendo mientras la corriente sea superior a la corriente de mantenimiento $I_H$.

% ============================================================
% 7. AISLAMIENTO GALVÁNICO
% ============================================================

\subsection{Necesidad de los tres aislamientos galvánicos}

El circuito utiliza tres mecanismos de aislamiento con funciones complementarias:

\begin{enumerate}
    \item \textbf{Transformador T1:} proporciona aislamiento entre la red y la etapa utilizada para detectar el cruce por cero.
    
    \item \textbf{MOC3020:} aísla el microcontrolador de la etapa de disparo y potencia del SCR.
    
    \item \textbf{IL300:} aísla la medición analógica de la tensión de carga respecto al microcontrolador.
\end{enumerate}

Estos aislamientos permiten intercambiar información entre las diferentes etapas sin establecer una conexión eléctrica directa.

La separación también reduce el riesgo de que transitorios de la etapa de potencia lleguen al microcontrolador, al computador o a los instrumentos de medición.

Por esta razón:

\[
GND_{\mathrm{control}}
\neq
Referencia_{\mathrm{potencia}}
\]

Mantener estas referencias separadas es fundamental para conservar el aislamiento galvánico del sistema.

% ============================================================
% 8. SCR VS TRIAC
% ============================================================

\subsection{Diferencia entre SCR y TRIAC}

El SCR es un dispositivo semiconductor unidireccional. Una vez disparado mediante su compuerta, conduce en un único sentido y permanece activado mientras la corriente principal no disminuya por debajo de la corriente de mantenimiento.

El TRIAC, en cambio, puede conducir corriente en ambos sentidos, por lo que resulta especialmente apropiado para controlar cargas AC directamente.

En el circuito de la práctica, el MOC3020 posee una salida tipo TRIAC, pero no controla directamente la potencia del bombillo. Su función es proporcionar aislamiento y transmitir el pulso necesario para activar el SCR.

% ============================================================
% 9. CONCLUSIONES
% ============================================================

\section{Conclusiones}

El laboratorio permitió comprobar experimentalmente el funcionamiento del control de potencia por ángulo de fase mediante el SCR BT151-500R. La variación del ángulo de disparo modificó directamente el intervalo de conducción del dispositivo y, en consecuencia, la tensión y potencia entregadas al bombillo.

La etapa de detección de cruce por cero permitió sincronizar el microcontrolador con la señal de la red. A partir de esta referencia, el microcontrolador generó el pulso de disparo con un retardo asociado al ángulo $\alpha$.

El transistor Q1 y el MOC3020 permitieron realizar el disparo del SCR sin conectar eléctricamente la etapa de control con la etapa de potencia. El aislamiento galvánico resultó fundamental para proteger el microcontrolador y reducir la posibilidad de transferencia de perturbaciones provenientes de la red.

La medición mediante el IL300 permitió obtener una representación aislada de la tensión de la carga. El LM358 acondicionó esta señal para que pudiera ser procesada por la entrada analógica A0 del microcontrolador.

Finalmente, se verificó que un pulso de Gate insuficiente o demasiado corto puede producir disparos inestables del SCR. Por tanto, el diseño del sistema de disparo debe considerar simultáneamente el ángulo de activación, la corriente de compuerta y la duración del pulso.

En conjunto, los resultados permiten concluir que la combinación de sincronización, control de fase, aislamiento galvánico y realimentación analógica constituye una estrategia efectiva para regular de manera controlada la potencia suministrada a una carga de corriente alterna.

\end{document}
